\section{Statistics and Linear Algebra}

\answer{AG1-1}
Since $\Cov(W,Z)=\Var(T)$, the latent variance is $3$. Then $\Var(U)=5-3=2$ and $\Var(V)=4-3=1$, both nonnegative. The instrumental-variable ratio gives $\beta=6/3=2$, whereas the ordinary slope is $6/5$. Independent centered normal variables with variances $3,2,1$, together with independent centered response noise, realize the specified moments.
The model necessarily imposes $\Cov(W,Y)=\Cov(Z,Y)=\beta\Var(T)$. Values $6$ and $5$ therefore contradict it; changing measurement-noise variances cannot repair the contradiction. This is a testable model restriction, not a reason to choose one ratio arbitrarily.

\answer{AG1-2}
A25 bounds both probabilities by $1/(1+2^2)=1/5$. The non-strict event attains it with probability $1/5$ at $2$ and $4/5$ at $-1/2$. The mean is $2/5-2/5=0$, and the second moment is $4/5+1/5=1$.
The strict event cannot attain the bound: equality would also force equality for the non-strict event, hence the distribution just found, whose strict-tail probability is zero. For $h>0$, let $b=2+h$ and put probability $1/(1+b^2)$ at $b$ and $b^2/(1+b^2)$ at $-1/b$. Its mean is zero and variance one. Its strict-tail probability tends to $1/5$ as $h\downarrow0$.

\answer{AG1-3}
Ignoring a constant, the objective is $\sum_j(b_j-z_j)^2/2+\sum_jb_j/2$. With budget multiplier $\nu\ge0$, coordinatewise minimization under nonnegativity gives
\[
 b_j=(z_j-1/2-\nu)_+.
\]
At $\nu=0$ this is $(5/2,1/2,0)$, whose sum is $3$, violating the budget. The budget must bind. Taking $\nu=1/2$ gives $(2,0,0)$ with sum $2$. The first coordinate has zero adjusted gradient; the second has zero right derivative at zero; the third has positive right derivative. All optimality conditions hold. The quadratic part is strictly convex, so the solution is unique even though the second coordinate is exactly at the threshold. The common multiplier couples the coordinates.

\answer{AG1-4}
Take $c=(1,1)^\top,d=0$. In A06, $e=1$ and $q=1+1/4=5/4$, so
\[
 \widehat\beta_c=(2,-1)^\top-\frac45(1,1/4)^\top
 =(6/5,-6/5)^\top,
 \qquad\Delta\mathrm{RSS}=4/5.
\]
For a direct check, put $\beta=(t,-t)$. The extra sum of squares is $(t-2)^2+4(1-t)^2$, whose derivative $10t-12$ vanishes uniquely at $t=6/5$. Substitution gives $16/25+4/25=4/5$. Equal-sized corrections to the two coefficients would therefore not be optimal.

\par\Needspace{8\baselineskip}
\subsection*{Section Exercises}
\answer{OA02}
Centering does not change slopes, so take all means to be zero. Write the joint regression as
\[
 Y=bX_1+\beta_2X_2+u,
 \qquad \Cov(u,X_1)=\Cov(u,X_2)=0.
\]
Set $\gamma=\Cov(X_1,X_2)/\Var(X_1)$ and let $r=X_2-\gamma X_1$ be the residual from regressing $X_2$ on $X_1$. The first normal equation implies that the simple-regression slope of $Y$ on $X_1$ is $b+\beta_2\gamma$. Thus
\[
 e=\beta_2r+u.
\]
Since $r$ is uncorrelated with $X_1$ and $X_2=r+\gamma X_1$,
\[
 \Cov(r,X_2)=\Var(r)
 =\Var(X_2)(1-\rho^2).
\]
Consequently,
\[
 \boxed{\beta_1
 =\frac{\Cov(e,X_2)}{\Var(X_2)}
 =(1-\rho^2)\beta_2.}
\]
Residualizing only the response leaves $\Var(X_2)$ in the denominator of the second slope. Regressing $e$ on $r$ instead uses $\Var(r)$ and recovers $\beta_2$. The two slopes agree when $\rho=0$. At $|\rho|=1$, the two joint-regression coefficients are generally not separately identifiable, which explains the restriction in the question.

\answer{OA03}
Write $A=\{X+Y>1\}$. The original joint density is $1$ on the unit square, and $\Prob(A)=1/2$, so the conditional density is $2$ on the upper triangle. Given $X=x$ and $A$, the variable $Y$ is uniform on $(1-x,1)$. Hence
\[
 \E[Y\mid X=x,A]=1-\frac x2,
 \qquad 0<x<1.
\]
The conditional mean is already affine. The residual $u=Y-(1-X/2)$ satisfies $\E[u\mid X,A]=0$, so it is orthogonal to both the constant and $X$. The regression therefore has intercept $1$ and slope
\[
 \boxed{-\frac12}.
\]
For a moment check, the conditional density $f_{X\mid A}(x)=2x$ gives
\[
 \E[X\mid A]=\frac23,\quad
 \E[X^2\mid A]=\frac12,\quad
 \Var(X\mid A)=\frac1{18}.
\]
The conditional-mean formula implies $\Cov(X,Y\mid A)=-\Var(X\mid A)/2=-1/36$. Dividing by the conditional variance again gives $-1/2$. Independence holds before selection and cannot be carried over to the conditional population.

\answer{OA05}
Let $c=X^\top y/n$. Expanding the objective difference gives
\[
 L(\beta)-L(0)
 =\frac{\|X\beta\|^2}{2n}-c^\top\beta
   +\lambda\|\beta\|_1.
\]
If $\|c\|_\infty\le\lambda$, the inequality $c^\top\beta\le\|c\|_\infty\|\beta\|_1$ makes this difference nonnegative, so zero is optimal.
Conversely, suppose $|c_j|>\lambda$ for some coordinate. Take $\beta=t\operatorname{sign}(c_j)e_j$ with $t>0$. Then
\[
 L(\beta)-L(0)
 =\frac{t^2\|X_j\|^2}{2n}-t(|c_j|-\lambda).
\]
For sufficiently small $t$, the negative linear term dominates the quadratic term. Thus the objective decreases, proving that the required condition is
\[
 \boxed{\lambda\ge\frac{\|X^\top y\|_\infty}{n}}.
\]
Under this condition, zero is also unique when $\lambda>0$. A second minimizer with $X\beta\ne0$ would, by strict convexity in the fitted value, have a midpoint with zero that improves the objective. A nonzero vector with $X\beta=0$ strictly increases the penalty. When $\lambda=0$, the condition becomes $X^\top y=0$, and the full minimizer set is $\ker X$; zero is then unique exactly when $X$ has full column rank.

\answer{OA30}
The relation is
\[
 \boxed{\lambda_i\le\mu_i\le\lambda_{i+1}},
 \qquad i=1,\ldots,n-1.
\]
Let $H=\{x\in\R^n:x_k=0\}$. Inserting a zero in coordinate $k$ identifies the Rayleigh quotient of $B$ with that of $A$ restricted to $H$. Write $R_A(x)=x^\top Ax/\|x\|^2$. The min--max principle gives
\[
 \lambda_i=\min_{\dim S=i}\ \max_{0\ne x\in S}R_A(x).
\]
The formula for $\mu_i$ minimizes the same quantity only over subspaces $S\subset H$. Restricting the choices cannot decrease the minimum, so $\mu_i\ge\lambda_i$.

For the other bound, let $E$ be spanned by eigenvectors of $A$ corresponding to its first $i+1$ eigenvalues. Since $H$ has codimension one, $\dim(E\cap H)\ge i$. Choose an $i$-dimensional subspace $S$ of this intersection. Every Rayleigh quotient on $S$ is at most $\lambda_{i+1}$. Using this candidate in the min--max formula for $\mu_i$ proves $\mu_i\le\lambda_{i+1}$.

Repeated eigenvalues are allowed, so strict inequalities cannot hold in general. For the identity matrix, equality holds on both sides.

